%%%%%%%%%%%%%%%%%%%%%%%%%%%%%%%%%%%%%%%%%%%%%%%%%%%%%%%%%%%%%%%%%%%%%%%%%%%%%%%%%
%	TASK 1 - Problem Definition and Experimental Design
%   Maps to brief Tasks 1.1, 1.2, 1.3.
%%%%%%%%%%%%%%%%%%%%%%%%%%%%%%%%%%%%%%%%%%%%%%%%%%%%%%%%%%%%%%%%%%%%%%%%%%%%%%%%%
\lhead{Task 1. \emph{Problem Definition and Experimental Design}}
\chapter{Task 1: Problem Definition and Experimental Design}
\label{Task1}

\section{Benchmark TSP Instance (Task 1.1)}
\label{T1:Benchmark}

We use the \textbf{kroA100} instance from the TSPLIB
collection~\cite{reinelt1991tsplib,applegate2006traveling}. It contains 100
cities in the Euclidean plane, uses symmetric \texttt{EUC\_2D} edge weights,
and has a known optimal tour length of 21{,}282. The number of cities falls within
the 70--150 range required by the project brief.

kroA100 is one of the most cited TSPLIB instances in the GA-for-TSP
literature, so its difficulty and optimal value are well known. At
$n=100$ the search space contains roughly $99!/2 \approx 4.7\times10^{155}$
distinct tours, which rules out exact enumeration. At the same time, the
instance is small enough for a 54-configuration parameter sweep with
30 runs each within a reasonable compute budget. Having a known optimum also
lets us report results as a relative gap rather than raw tour length.

\begin{table}[htbp]
	\centering
	\caption{kroA100 benchmark instance metadata.}
	\label{tab:instance}
	\begin{tabular}{ll}
		\hline
		\textbf{Property} & \textbf{Value} \\
		\hline
		Name & kroA100 \\
		Number of cities ($n$) & 100 \\
		Edge-weight type & \texttt{EUC\_2D} (symmetric Euclidean 2D) \\
		Known optimal tour length & 21{,}282 \\
		Instance class & Symmetric Euclidean TSP \\
		Source & TSPLIB~\cite{reinelt1991tsplib,applegate2006traveling} \\
		\hline
	\end{tabular}
\end{table}

For \texttt{EUC\_2D} instances, each city $i$ has coordinates $(x_i, y_i)$ and
the inter-city distance follows the TSPLIB convention of nearest-integer
rounding,
\[
d(i,j) = \left\lfloor \sqrt{(x_i - x_j)^2 + (y_i - y_j)^2} + 0.5 \right\rfloor
\]
so that $d(i,j) = d(j,i)$. We use the same rounded distance for fitness
evaluation and for comparison against the published optimum.

\section{Formal Mathematical Formulation (Task 1.2)}
\label{T1:Formulation}

We formulate the symmetric TSP on kroA100 directly as a permutation
optimisation problem, matching the GA representation used in the rest of the
report. Let $V = \{0, 1, \ldots, n-1\}$ denote the set of city indices with
$n=100$, and let $d : V \times V \to \mathbb{Z}_{\geq 0}$ denote the rounded
Euclidean distance defined above.

A candidate solution is a tour $\pi = (\pi_0, \pi_1, \ldots, \pi_{n-1})$,
where $\pi$ is a permutation of $V$. The fitness function to be minimised is
the total length of the closed tour induced by $\pi$,
\[
\min_{\pi} \; f(\pi)
\;=\;
\sum_{i=0}^{n-1} d\!\left(\pi_i,\, \pi_{(i+1) \bmod n}\right),
\]
where the modular index in the final term closes the tour by linking
$\pi_{n-1}$ back to $\pi_0$.

Two constraints are implicit in this formulation. The permutation constraint
requires that every city is visited exactly once, i.e.\
$\{\pi_0, \pi_1, \ldots, \pi_{n-1}\} = V$ with no repetitions. The
closed-tour constraint is captured by the modular indexing and forces the
salesman to return to the starting city. Together they replace the
incoming-degree, outgoing-degree, and subtour-elimination constraints of the
classical arc-based integer programming formulation. Any permutation of $V$
is automatically a Hamiltonian cycle, so feasibility comes down to keeping
the permutation property intact under the GA operators.

The search space is discrete and finite. After removing rotational and
reflective symmetries, the number of distinct tours is $\dfrac{n!}{2n}$, which
for $n=100$ is approximately $4.7 \times 10^{155}$. This factorial growth
rules out exact search and motivates a population-based metaheuristic such as
a Genetic Algorithm~\cite{holland1975adaptation,goldberg1989genetic}.
The permutation constraint also affects operator design: standard
string-style crossover and bit-flip mutation do not preserve permutations,
so we either use feasibility-preserving operators or pair an unconstrained
operator with a repair step. We choose the second option so that repair can
be studied as an experimental factor, following the constraint-handling
taxonomy in Larra{\~n}aga et al.~\cite{larranaga1999genetic}.

\section{Experimental-Setup Justification (Task 1.3)}
\label{T1:Setup}

We set the experimental protocol before collecting any results so that every
configuration is evaluated under the same conditions.

Each configuration is run for \textbf{30 independent random seeds}. Thirty
seeds provide a reasonable sample for reporting means and standard deviations
while keeping the full sweep manageable.

Each run terminates at a fixed budget of \textbf{500 generations}.
Preliminary convergence curves (\sref{T4:Convergence}) show that most fitness
improvement happens within roughly the first 200 generations; the remaining
300 generations cover the slower phase where late mutations or rare crossover
events still shorten a few edges. A fixed generation budget, rather than an adaptive stopping
rule, ensures that every run produces the same number of per-generation
observations, keeping convergence plots, diversity curves, and statistical
tests directly comparable.

The parameter sweep covers four GA factors, giving
$3 \times 3 \times 3 \times 2 = 54$ unique configurations:

\begin{align*}
	\text{pop\_size} &\in \{50, 100, 200\} \\
	p_c &\in \{0.70, 0.85, 0.95\} \\
	p_m &\in \{0.01, 0.05, 0.10\} \\
	\text{selection} &\in \{\text{tournament},\, \text{roulette}\}
\end{align*}

Population size covers the moderate range commonly used in GA-TSP studies on
TSPLIB-scale instances. Crossover and mutation rates span the
low/medium/high settings from the empirical reviews of Larra{\~n}aga et
al.~\cite{larranaga1999genetic}, so the sweep can separate exploration
pressure from disruption. We include selection method at two levels because
tournament and roulette selection create different pressure
profiles. The interaction between selection pressure and mutation rate is
one of the questions \cref{Task4} addresses.

Combining the sweep with 30 seeds gives $54 \times 30 = 1{,}620$ total runs,
which is enough for non-parametric tests and a per-factor sensitivity
analysis.

A second group of parameters is kept fixed across the sweep because we want
to isolate the four factors above instead of exploring the full operator space.
Tournament size is $k=3$, elitism carries over 2 individuals per generation,
crossover uses the deliberately naive single-point operator with repair
enabled, and mutation uses the swap operator. Using naive single-point
crossover with repair is itself an experimental decision: it ensures that
every offspring needs repair, which is what \cref{Task2} and
\cref{Task3} require for the constraint-handling comparison.

Every run is identified by the pair (configuration hash, random seed). The
configuration hash is a digest of the full parameter dictionary. Per-generation
best fitness, mean fitness, and Hamming diversity are written to CSV files
keyed by that pair, and final tours are stored separately. This lets any
reported figure, table, or statistical test be regenerated from the stored
data without re-running the GA.
